%قبل از شروع مسئله ترجیحاً چیزی ننویسید. اما اگر خواستید بنویسید، در قالب لیستی (فقرات و شمارش) بنویسید.
%برای اضافه کردن نمره، بعدی از شروع‌مسئله، [x] را اضافه کنید و به جای x  نمره مورد نظر را قرار دهید. برای زیرمسئله‌ها هم به همین صورت است.
%\شروع‌مسئله[10]
\شروع‌مسئله [12]
%\متن‌سیاه{سوال عملی اول:}\\
%\متن‌سیاه{دراپ‌تست}\\
یکی از وب‌سایت‌های اخبار تکنولوژی پس از معرفی آیفون جدید قصد دارد برای جذب مخاطب اقدام به برگزاری یک آزمون پرتاب\پانویس{drop test} مهیج کند. آن‌ها یک ساختمان $k$ طبقه ($1\leq k\leq 10^4$) در اختیار دارند و می‌خواهند بدانند آیفون جدید سقوط از حداکثر کدام طبقه را می‌تواند تحمل کند. از آنجا که کوتاه بودن ویدیوی نهایی و سر نرفتن حوصلۀ کاربران برای این وب‌سایت مهم است، آن‌ها می‌خواهند آزمون سقوط در پرتاب‌های کمتری خاتمه یابد. همچنین با توجه به افزایش قیمت دلار، این وب‌سایت می‌خواهد از حداکثر دو گوشی آیفون در این آزمون‌ها استفاده کند.
فرض کنید در نتیجۀ هر آزمایش یا گوشی کاملا سالم می‌ماند یا نابود می‌شود. تنها در حالتی که گوشی سالم بماند می‌توان از آن برای آزمایش‌های بعدی نیز استفاده کرد.
آن‌ها به سراغ شما آمده‌اند تا به آن‌ها کمک کنید.

در بستۀ تمرین، یک فایل \متن‌تایپ{droptest.py} در اختیار شما قرار داده شده که دارای یک متد \متن‌تایپ{test} است که باید آن را بهینه کنید. این متد دو ورودی دارد: یک شیء \متن‌تایپ{iphone} و یک عدد $k$ (تعداد طبقات ساختمان). شما می‌توانید از شیء \متن‌تایپ{iphone}، متد \متن‌تایپ{drop(floor)} را اجرا کنید که اگر در نتیجۀ آزمونِ پرتاب از طبقۀ \متن‌تایپ{floor}، گوشی سالم بماند ۰ و اگر بشکند ۱ برمی‌گرداند. در حالت‌های دیگر (مثل وقتی که \متن‌تایپ{floor} عدد معتبری نباشد یا دو گوشی قبلاً شکسته شده باشند)، این تابع $-1$ برمی‌گرداند. متد \متن‌تایپ{test} در نهایت باید پایین‌ترین طبقه‌ای که پرتاب گوشی از آن منجر به شکستن می‌شود را برگرداند و اگر چنین طبقه‌ای وجود ندارد، $k+1$ برگرداند.

متد نمونه موجود در بستۀ تمرین، دارای یک الگوریتم بدیهی است که طبقه‌ها را از پایین به بالا تست می‌کند و اولین طبقه‌ای که با پرتاب گوشی از آن طبقه، گوشی بشکند را برمی‌گرداند. شما باید این متد را بهینه کنید و تعداد آزمون‌ها را کمتر کنید. قبل از ارسال برنامۀ خود، آن را با استفاده از دستور \متن‌تایپ{python sampletester.py} تست کنید. این تست برنامۀ شما را در حالتی اجرا می‌کند که یک ساختمان ۱۰۰ طبقه وجود دارد که از طبقۀ ۹۹اُم به بالا اگر گوشی را پرت کنیم می‌شکند. در نهایت فایل \متن‌تایپ{droptest.py} جدید را در کوئرا ارسال کنید. نمرۀ شما بسته به این دارد که برنامۀ شما از چند آزمون استفاده می‌کند. 

%\متن‌سیاه{ورودی:}\\
%در خط اول ورودی، ابتدا k و سپس m داده می‌شود.\\
%\متن‌سیاه{خروجی:}\\
%در تنها سطر خروجی، تعداد آیفون‌های مورد نیاز می‌آید. اگر در هیچ صورتی چنین کاری ممکن نبود ۱- چاپ می‌شود.\\
%\متن‌سیاه{محدودیت‌ها:}\\
%$1\leq k\leq 10^4$

%$m\leq10^9$
%\\
%\textbf{\متن‌سیاه{ورودی و خروجی نمونه}\\
%% تست اول
%\begin{latin}
%	\begin{tabularx}{\textwidth}{|X | X|}
%		\hline
%		\textbf{stdin} &
%		\textbf{stdout}\\
%		\hline
%		100 100&
%		1\\
%		\hline
%	
%	\end{tabularx}
%\end{latin}
%
%% تست دوم
%
%\begin{latin}
%	\begin{tabularx}{\textwidth}{|X | X|}
%		\hline
%		\textbf{stdin} &
%		\textbf{stdout}\\
%		\hline
%		340 20&
%		3\\
%		\hline
%		
%	\end{tabularx}
%\end{latin}}
% صورت اصلی مسئله را در این‌جا بنویسید.
%اگر مسئله چند قسمت داشت به این صورت عمل کنید:
%\شروع{زیرمسائل}
%\زیرمسئله[0] % مسئلۀ ۱آ
%\زیرمسئله % مسئلۀ ۱ب
%\زیرمسئله % مسئلۀ ۱پ
%\زیرمسئله % مسئلۀ ۱ت
%\پایان{زیرمسائل}

\solution{
	پاسخ نمونه.
}




%%%%%%%%%%%%%%%%%%%%%%%%%%%%%%%%%%%%%%%%%%%%%%%%%%%%%
% برخی دستورات مفید برای تمرین‌ها. %
%%%%%%%%%%%%%%%%%%%%%%%%%%%%%%%%%%%%%%%%%%%%%%%%%%%%%

% فرمول‌های نمونه:
%$f(x) = \Theta(x)$
%$T(x, y) \leq \log(x) + 2^y + \binom{2n}{n}$
% {\tt code\_function}


% لیست موارد بدون شماره‌گذاری:
%\شروع{فقرات}
%    \فقره مورد اول 
%        \شروع{فقرات}
%            \فقره مورد اول داخلی 
%                \شروع{فقرات}
%                    \فقره اولین مورد خیلی داخلی 
%                    \فقره دومین مورد خیلی داخلی 
%                \پایان{فقرات}
%            \فقره مورد دوم داخلی 
%        \پایان{فقرات}
%    \فقره مورد دوم 
%\پایان{فقرات}

% لیست موارد شماره‌گذاری شده:
%\شروع{شمارش}
%    \فقره یک مورد 
%    \فقره یک مورد دیگر 
%    \فقره باز هم یک مورد دیگر
%\پایان{شمارش}

% برای نوشتن یک فرمول چند خطی که یک شماره بخوره. برای حالت‌های دیگر به سایت زیر مراجعه کنید:
% https://www.overleaf.com/learn/latex/Aligning_equations_with_amsmath
%\begin{align} 
%    \begin{split}
%        (x+y)^3 &= (x+y)^2(x+y) \\
%                &= (x^2+2xy+y^2)(x+y) \\
%                &= (x^3+2x^2y+xy^2) + (x^2y+2xy^2+y^3) \\
%                &= x^3+3x^2y+3xy^2+y^3
%    \end{split}                                 
%\end{align}

% نوشتن ماتریس و آرایه دوبعدی
%\begin{align}
%    A = 
%    \begin{bmatrix}
%        A_{11} & A_{21} \\
%        A_{21} & A_{22}
%    \end{bmatrix}
%\end{align}

% اضافه کردن تصویر با شرح. برای اضافه کردن تصویر بدون شرح، فقط دستور درج‌تصویر را استفاده کنید.
%\شروع{شکل}[h!]
%\تنظیم‌ازوسط
%  \درج‌تصویر[پهنا=0.5\پهنای‌متن]{img.jpg}
%\شرح{شرح تصویر}
%\برچسب{تصویر:۱}
%\پایان{شکل}
%\پایان{زیرمسائل}

%کد زیر را ببینید.
%%\lstinputlisting{source_filename.py}
%\begin{latin}
%\begin{lstlisting}
%# Import the modules
%import sys
%import random
%
%ans = True
%
%while ans:
%question = raw_input("Ask the magic 8 ball a question: (press enter to quit) ")
%
%answers = random.randint(1,8)
%
%if question == "":
%sys.exit()
%
%elif answers == 1:
%print "It is certain"
%
%elif answers == 2:
%print "Outlook good"
%
%elif answers == 3:
%print "You may rely on it"
%
%elif answers == 4:
%print "Ask again later"
%
%elif answers == 5:
%print "Concentrate and ask again"
%
%elif answers == 6:
%print "Reply hazy, try again"
%
%elif answers == 7:
%print "My reply is no"
%
%elif answers == 8:
%print "My sources say no"
%\end{lstlisting}
%\end{latin}
